% main.tex — TAGA-LLM: ICGNC 2026
% Springer svproc template
\documentclass{styles/svproc}

% --- 推荐宏包 ---
\usepackage{url}
\def\UrlFont{\rmfamily}
\usepackage{marvosym}
\usepackage{amsmath}
\usepackage{amssymb}
\usepackage{graphicx}
\graphicspath{{figures/}}
\usepackage{booktabs}
\usepackage{siunitx}
\usepackage{algorithm}
\usepackage{algorithmic}
\usepackage{multirow}
\usepackage{bm}
\usepackage{paralist}   % inparaenum 行内列表
\usepackage[colorlinks=true,linkcolor=blue,citecolor=blue,urlcolor=blue]{hyperref}

% 自定义命令
\DeclareMathOperator{\concat}{\|}  % 拼接符号

\begin{document}
\mainmatter

% ============================
% 标题
% ============================
\title{TAGA-LLM: Topology-Adaptive Graph Attention Informed\\
Large Language Model for Multi-UAV Cooperative\\
Search and Attack Mission Planning}

\titlerunning{TAGA-LLM for Multi-UAV Mission Planning}

\author{Xinrui Guo\inst{1,2} \and Jiang Zhao\inst{3} \and Yingxun Wang\inst{4}
\and Pei Chi\inst{4}}

\authorrunning{Xinrui Guo et al.}

\tocauthor{Xinrui Guo, Jiang Zhao, Yingxun Wang, Pei Chi}

\institute{
School of Artificial Intelligence, Beihang University, Beijing 100191, China
\and
Zhongguancun Academy, Beijing 100191, China
\and
School of Automation Science and Electrical Engineering, Beihang University, Beijing 100191, China
\and
Institute of Unmanned System, Beihang University, Beijing 100191, China
\\
\email{XinruiGuo@buaa.edu.cn}
}

\maketitle

% ============================
% 摘要
% ============================
\begin{abstract}
Cooperative search and attack missions for unmanned aerial vehicle (UAV) swarms
require real-time coordination under dynamic communication topologies, yet
existing approaches either rely on hand-crafted heuristics that lack
adaptability or employ multi-agent reinforcement learning (MARL) policies that
offer no interpretability.  We propose TAGA-LLM, a two-stage framework
that first trains a Topology-Adaptive Graph Attention enhanced Multi-Agent PPO
(TAGA-MAPPO) expert policy, then distills its decision-making capability into a
Qwen-4B large language model via LoRA fine-tuning with chain-of-thought
supervision.  The TAGA mechanism augments standard graph attention with three
learnable, physics-informed factors---distance decay, task awareness, and
information freshness---to model realistic communication constraints. Through
expert trajectory collection, counterfactual data augmentation, and a
topology-aware prompt encoding scheme, the distilled LLM achieves
instruction-controllable and interpretable mission planning while preserving
cooperative performance.  We evaluate the framework on a $100\!\times\!100$ grid
battlefield with 10~UAVs, 5~targets, and 7~threat zones, assessing both
cooperative task performance and the quality of chain-of-thought reasoning
across five operational directives.  A short CPU sanity run (random versus
greedy TAGA-MAPPO after $90$~PPO updates) confirms stable optimization under
curriculum scheduling and improved search coverage; full-scale statistical
comparisons are deferred to the extended evaluation campaign.
\keywords{multi-UAV cooperation, graph attention network, large language model,
multi-agent reinforcement learning, knowledge distillation}
\end{abstract}

% ============================
% 正文章节
% ============================
% sec_introduction.tex — Section 1: Introduction (compressed)
% ICGNC 2026, Springer svproc template

\section{Introduction}
\label{sec:intro}

Multi-UAV cooperative search and attack missions require a swarm of UAVs to
autonomously explore unknown environments, discover hidden targets, and
coordinate strikes under stringent communication constraints and threat
zones~\cite{shakhatreh2019unmanned}.  The core challenge is achieving effective
inter-agent coordination when each UAV possesses only a partial, local view and
can communicate only within a limited radius.

Multi-Agent Reinforcement Learning (MARL) offers a data-driven approach to
learning cooperative policies~\cite{lowe2017multi,yu2022surprising}.
Graph-based communication mechanisms such as
CommNet~\cite{sukhbaatar2016learning} and DGN~\cite{jiang2020graph} enable
learned message passing, yet standard Graph Attention Networks
(GATs)~\cite{velickovic2018gat} treat all links equally, ignoring that wireless
channel quality degrades with distance, co-modal agents share higher mutual
information, and stale messages carry increasing uncertainty.  Furthermore, MARL
policies are inherently opaque.  Large Language Models
(LLMs)~\cite{brown2020language} have shown strong reasoning
capabilities~\cite{yao2023react,shinn2023reflexion}, but directly applying them
to multi-UAV coordination is hindered by the lack of domain-specific
cooperative strategies and the difficulty of encoding topology-dependent states
into prompts.

We propose TAGA-LLM, a two-stage framework that combines MARL's
cooperative learning strength with LLM interpretability.  The key insight is
that MARL discovers effective cooperative strategies but locks them inside
opaque neural networks, while LLMs excel at structured reasoning and natural
language explanation but lack domain-specific cooperative knowledge;
distillation bridges this complementary gap.  In \emph{Stage~I}, we introduce the
\emph{Topology-Adaptive Graph Attention} (TAGA) mechanism with three learnable
factors---distance decay, task awareness, and information freshness---integrated
into a hierarchical MAPPO framework with curriculum learning.  In
\emph{Stage~II}, we distill the TAGA-MAPPO policy into a Qwen-4B LLM via LoRA
fine-tuning with topology-aware prompt encoding, CoT supervision, and
counterfactual data augmentation.

The primary contributions of this paper are summarized as follows:
\begin{inparaenum}[(a)]
\item a TAGA mechanism with three physics-informed adaptive factors for
  communication aware multi-agent coordination;
\item a two-stage MARL to LLM distillation pipeline that yields
  interpretable and instruction-controllable mission planning; and
\item a hierarchical policy architecture trained with MAPPO and a
  three-stage curriculum.
\end{inparaenum}

The rest of this paper is organized as follows.
In Sect.~\ref{sec:related}, we review related work.
In Sect.~\ref{sec:problem}, we state the Dec-POMDP formulation of the mission.
In Sect.~\ref{sec:method}, we present the TAGA-LLM framework, including TAGA,
MAPPO training, and distillation.
In Sect.~\ref{sec:experiments}, we report simulation and evaluation plans.
Finally, Sect.~\ref{sec:conclusion} concludes the paper.

% sec_related.tex — Section 2: Related Work (compressed)
% ICGNC 2026, Springer svproc template

\section{Related Work}
\label{sec:related}

We briefly review multi-agent communication learning, LLM-based decision
systems, policy distillation, and multi-UAV planning, and position TAGA-LLM
relative to these lines of work.

\paragraph{Communication in MARL.}
CommNet~\cite{sukhbaatar2016learning} broadcasts mean-pooled messages;
TarMAC~\cite{das2019tarmac} adds attention-based targeting;
DGN~\cite{jiang2020graph} uses graph convolutions but assumes a static
adjacency matrix.  QMIX~\cite{rashid2018qmix} and
MAPPO~\cite{yu2022surprising} address credit assignment without modeling
physical communication constraints.  Our TAGA mechanism injects three
physics-informed factors directly into the attention computation, treating the
graph as a dynamic, quality-variable channel.

\paragraph{LLMs for decision-making.}
ReAct~\cite{yao2023react} interleaves reasoning with actions;
Reflexion~\cite{shinn2023reflexion} adds self-reflection;
Voyager~\cite{wang2023voyager} generates executable code;
SayCan~\cite{ahn2022saycan} grounds LLM plans in affordances.  These
approaches either operate in single-agent settings or rely on pre-defined
protocols, and none encodes topology-dependent multi-agent states into LLM
prompts while preserving MARL-learned cooperative strategies.

\paragraph{Knowledge distillation for policy transfer.}
Policy distillation compresses a teacher policy into a smaller
student~\cite{rusu2016policy}.  Recent efforts explore distilling RL-learned
behaviors into language models for
interpretability~\cite{brown2020language,li2023camel}, yet existing pipelines
do not account for graph-structured multi-agent communication or
topology-dependent state encoding.  Our Stage~II distillation combines
LoRA-based parameter efficient fine-tuning~\cite{hu2022lora} with
topology-aware prompt encoding and chain-of-thought supervision to preserve
cooperative strategies in the distilled LLM.

\paragraph{Multi-UAV planning.}
Classical methods include potential fields~\cite{ge2002dynamic}, auction-based
allocation~\cite{choi2009consensus}, and metaheuristics~\cite{dorigo2006ant}.
MARL-based approaches such as MADDPG~\cite{lowe2017multi} and
HAPPO~\cite{kuba2022trust} train end-to-end policies, but for homogeneous
swarms, parameter-sharing MAPPO is more efficient~\cite{yu2022surprising}.  Our
work builds on MAPPO with TAGA for communication-aware coordination and bridges
the MARL--LLM gap through policy distillation.

% sec_problem.tex — Section 3: Problem Statement (ICGNC proceedings style)
% ICGNC 2026, Springer svproc template

\section{Problem Statement}
\label{sec:problem}

We model the mission as a Decentralized Partially Observable Markov Decision
Process (Dec-POMDP) $\langle \mathcal{N}, \mathcal{S}, \{\mathcal{O}_i\}_{i=1}^{N},
\{\mathcal{A}_i\}_{i=1}^{N}, \mathcal{T}, \mathcal{R}, \gamma \rangle$ with
$N\!=\!10$ homogeneous UAVs on a $100\!\times\!100$ grid containing $M\!=\!5$
hidden targets and $K\!=\!7$ threat zones.

\paragraph{Observation.}
Each UAV~$i$ observes $o_i = (o_i^{\text{grid}}, o_i^{\text{self}})$: a
five-channel $11\!\times\!11$ binary map (obstacles, explored cells, friendly
UAVs, targets, threats) and a 6-dim state vector (position, velocity, energy,
mode).

\paragraph{Action.}
UAV~$i$ selects mode $m_i \!\in\! \{\text{search}, \text{attack}\}$, then a
mode-conditioned action: 9 moves (8 directions + stay) in search mode, plus
\emph{fire} in attack mode.  Invalid actions are masked.

\paragraph{Communication.}
A dynamic graph $\mathcal{G}_t\!=\!(\mathcal{N}, \mathcal{E}_t)$ with edge
$(i,j) \!\in\! \mathcal{E}_t$ iff $d_{ij} \!\leq\! R_{\text{comm}}\!=\!20$.
For each pair $(i,j)$, we track the \emph{message staleness}
$\Delta t_{ij} \!\geq\! 0$, defined as the number of time steps since UAV~$j$
last successfully transmitted to UAV~$i$ (reset to 0 upon reception).

\paragraph{Reward.}
The shared reward combines search coverage ($w_s\!=\!0.1$), target discovery
($w_d\!=\!1.0$), destruction ($w_a\!=\!5.0$), cooperative bonus
($w_c\!=\!3.0$), threat/loss penalty ($w_p\!=\!{-}2.0$), and time cost
($w_t\!=\!{-}0.01$):
\begin{equation}
  r_t = w_s r_t^{\text{cov}} + w_d r_t^{\text{disc}} + w_a r_t^{\text{atk}}
  + w_c r_t^{\text{coop}} + w_p r_t^{\text{pen}} + w_t.
  \label{eq:reward}
\end{equation}
The terms $r_t^{\text{cov}},\ldots,r_t^{\text{pen}}$ denote incremental
contributions at step~$t$ and are typically sparse (zero unless the corresponding
event occurs).
We evaluate via five metrics: search coverage, target destroy rate,
completion time, UAV survival rate, and cooperative attack ratio.

% sec_method.tex — Section 4: Proposed Method (compressed)
% ICGNC 2026, Springer svproc template

\section{Proposed Method}
\label{sec:method}

This section details the TAGA-LLM framework.  As shown in
Fig.~\ref{fig:framework}, Stage~I (Sects.~\ref{sec:taga}--\ref{sec:mappo})
trains a TAGA-MAPPO expert policy, and Stage~II (Sect.~\ref{sec:distill})
distills it into an instruction-controllable LLM.

\begin{figure}[t]
\centering
\includegraphics[width=\textwidth]{figures/Figure1.png}
\caption{The TAGA-LLM framework.  Stage~I: TAGA-MAPPO with topology-adaptive
graph attention and curriculum learning.  Stage~II: distillation to Qwen-4B via
LoRA fine-tuning with chain-of-thought supervision.}
\label{fig:framework}
\end{figure}

% --- 4.1 TAGA Mechanism ---
\subsection{Topology-Adaptive Graph Attention}
\label{sec:taga}

Given node features $\mathbf{h}_i^{(0)} \!\in\! \mathbb{R}^{128}$ from the
observation encoder (Sect.~\ref{sec:encoder_policy}), the attention logit
between UAV~$i$ and neighbor~$j \!\in\! \mathcal{N}_i$ at head~$k$ is
\begin{equation}
  e_{ij}^{(k)} = \frac{(\mathbf{W}_Q^{(k)} \mathbf{h}_i^{(0)})^\top
  (\mathbf{W}_K^{(k)} \mathbf{h}_j^{(0)})}{\sqrt{d_k}},
  \label{eq:attn_logit}
\end{equation}
where $\mathbf{W}_Q^{(k)}, \mathbf{W}_K^{(k)} \!\in\! \mathbb{R}^{d_k \times 128}$
project into the $k$-th head space.
We form a \emph{modulated logit} by adding the log-domain TAGA biases to
$e_{ij}^{(k)}$, so that the three factors act as multiplicative weights on the
resulting attention coefficients:
\begin{equation}
  s_{ij}^{(k)} = e_{ij}^{(k)}
    + \log f_d(d_{ij})
    + \log f_m(m_i, m_j)
    + \log f_t(\Delta t_{ij}),
  \label{eq:taga_s}
\end{equation}
and normalize over neighbors with a standard softmax:
\begin{equation}
  \alpha_{ij}^{\text{TAGA},(k)} = \operatorname{softmax}_{j \in \mathcal{N}_i}\!\big(s_{ij}^{(k)}\big),
  \label{eq:taga_softmax}
\end{equation}
equivalently:
\begin{equation}
  \alpha_{ij}^{\text{TAGA},(k)} = \frac{
    \exp\!\big(e_{ij}^{(k)}\big)\, f_d\, f_m\, f_t
  }{
    \sum_{j' \in \mathcal{N}_i} \exp\!\big(e_{ij'}^{(k)}\big)\, f_d\, f_m\, f_t
  },
  \label{eq:taga}
\end{equation}
where $\mathcal{N}_i = \{j : (i,j) \!\in\! \mathcal{E}_t\}$ denotes the
communication neighbors of UAV~$i$ (only terms with $j \!\in\! \mathcal{N}_i$
appear in \eqref{eq:taga}), and the three modulation factors are:
(i)~\emph{distance decay}
$f_d(d_{ij}) = \exp(-\lambda_d \cdot d_{ij}/R_{\text{comm}})$, modeling
signal attenuation with distance (the indicator
$\mathbb{1}(d_{ij} \!\leq\! R_{\text{comm}})$ is implicit in the edge set
$\mathcal{E}_t$);
(ii)~\emph{task awareness}
$f_m(m_i, m_j) = 1 + \lambda_m \cdot \mathbb{1}(m_i \!=\! m_j)$, amplifying
attention between co-modal agents; (iii)~\emph{information freshness}
$f_t(\Delta t_{ij}) = \exp(-\lambda_t \cdot \Delta t_{ij})$, discounting
messages by their staleness $\Delta t_{ij}$ (cf.\ Sect.~\ref{sec:problem}).
The coefficients $\lambda_d, \lambda_m, \lambda_t \!\in\! \mathbb{R}_{>0}$
are learnable scalars (initialized to $1.0, 1.0, 0.1$, respectively).
The multi-head aggregated output is:
\begin{equation}
  \tilde{\mathbf{h}}_i = \operatorname{LN}\!\bigg(\mathbf{W}_O
  \concat_{k=1}^{H}\Big(\sum_{j \in \mathcal{N}_i}
  \alpha_{ij}^{\text{TAGA},(k)} \mathbf{W}_V^{(k)} \mathbf{h}_j^{(0)}\Big)
  + \mathbf{W}_R \mathbf{h}_i^{(0)}\bigg),
  \label{eq:taga_agg}
\end{equation}
where $\concat$ denotes concatenation across $H$ heads and $\mathbf{W}_R$ is a
residual projection handling dimension mismatch.  We stack $L\!=\!2$ such
layers ($H\!=\!4$ heads, $d_k\!=\!16$, output dim~64); the final-layer output
is denoted $\mathbf{h}_i^{(L)} \!\equiv\! \tilde{\mathbf{h}}_i
\!\in\! \mathbb{R}^{64}$.

% --- 4.2 Observation Encoder and Policy ---
\subsection{Observation Encoder and Hierarchical Policy}
\label{sec:encoder_policy}

The observation encoder has two branches:
\begin{equation}
  \mathbf{c}_i = \phi_{\text{cnn}}(\mathbf{x}_i) \in \mathbb{R}^{128},\qquad
  \mathbf{s}_i = \phi_{\text{mlp}}(\mathbf{u}_i) \in \mathbb{R}^{64},
  \label{eq:enc_branch}
\end{equation}
where $\phi_{\text{cnn}}$
($\operatorname{Conv}_{5\to32}^{5\times5} \!\to\!
\operatorname{Conv}_{32\to64}^{5\times5} \!\to\!
\operatorname{AvgPool}$) encodes the
$5\!\times\!11\!\times\!11$ local grid $\mathbf{x}_i$, and
$\phi_{\text{mlp}}$ ($6 \!\to\! 32 \!\to\! 64$) encodes the self-state
vector $\mathbf{u}_i$.  The concatenated features are linearly projected:
\begin{equation}
  \mathbf{h}_i^{(0)} = \phi_o([\mathbf{c}_i;\, \mathbf{s}_i])
  = \mathbf{W}_E [\mathbf{c}_i;\, \mathbf{s}_i] + \mathbf{b}_E
  \;\in\; \mathbb{R}^{128},
  \label{eq:enc_proj}
\end{equation}
which serves as input to the TAGA layers.

Let $\tilde{\mathbf{h}}_i = \mathbf{h}_i^{(L)} \!\in\! \mathbb{R}^{64}$
denote the output of the second TAGA layer.  We form
\begin{equation}
  \mathbf{f}_i = [\tilde{\mathbf{h}}_i;\, \mathbf{h}_i^{(0)}]
  \;\in\; \mathbb{R}^{192}.
  \label{eq:feat_concat}
\end{equation}
The mode, search, and attack heads are implemented as MLPs with one hidden
layer of width~128:
\begin{equation}
  \boldsymbol{\ell}_i^{m} = g_m(\mathbf{f}_i) \in \mathbb{R}^{2},\quad
  \boldsymbol{\ell}_i^{s} = g_s(\mathbf{f}_i) \in \mathbb{R}^{9},\quad
  \boldsymbol{\ell}_i^{a} = g_a(\mathbf{f}_i) \in \mathbb{R}^{10}.
  \label{eq:logits}
\end{equation}
Let $m_i \!\in\! \{\text{search},\, \text{attack}\}$ and
$a_i \!\in\! \mathcal{A}_{m_i}$, where
$|\mathcal{A}_{\text{search}}| \!=\! 9$ (8 directions + stay) and
$|\mathcal{A}_{\text{attack}}| \!=\! 10$ (9 moves + fire).
The hierarchical policy factorizes as
\begin{equation}
  \pi(m_i, a_i \mid \mathbf{f}_i)
  = \pi^{m}(m_i \mid \mathbf{f}_i)\;
    \pi(a_i \mid m_i, \mathbf{f}_i),
  \label{eq:policy}
\end{equation}
with
\begin{align}
  \pi^{m}(m_i \mid \mathbf{f}_i)
    &= \operatorname{softmax}(\boldsymbol{\ell}_i^{m}), \label{eq:pi_mode}\\
  \pi(a_i \mid m_i\!=\!\text{search},\, \mathbf{f}_i)
    &= \operatorname{softmax}\!\big(\operatorname{mask}_s(\boldsymbol{\ell}_i^{s})\big),
    \label{eq:pi_search}\\
  \pi(a_i \mid m_i\!=\!\text{attack},\, \mathbf{f}_i)
    &= \operatorname{softmax}\!\big(\operatorname{mask}_a(\boldsymbol{\ell}_i^{a})\big),
    \label{eq:pi_attack}
\end{align}
where $\operatorname{mask}_s$ and $\operatorname{mask}_a$ set invalid-action
logits to $-\infty$ before normalization.

% --- 4.3 MAPPO Training ---
\subsection{MAPPO Training with Curriculum Learning}
\label{sec:mappo}

We train with Centralized Training with Decentralized Execution (CTDE).
A centralized critic receives mean-pooled TAGA features and outputs a scalar
state value:
\begin{equation}
  \bar{\mathbf{h}}_t = \frac{1}{N_t}\sum_{i=1}^{N_t}
  \tilde{\mathbf{h}}_{i,t},\qquad
  V_\phi(\bar{\mathbf{h}}_t) \in \mathbb{R}.
  \label{eq:critic}
\end{equation}
Let $u_{i,t} = (m_{i,t},\, a_{i,t})$ denote the hierarchical action of
agent~$i$ at time~$t$.  Its probability factorizes as
\begin{equation}
  \pi_\theta(u_{i,t} \mid \mathbf{f}_{i,t})
  = \pi_\theta^{m}(m_{i,t} \mid \mathbf{f}_{i,t})\;
    \pi_\theta(a_{i,t} \mid m_{i,t},\, \mathbf{f}_{i,t}).
  \label{eq:mappo_pi}
\end{equation}
For shared-parameter MAPPO, the importance ratio is
\begin{equation}
  r_{i,t}(\theta) = \frac{
    \pi_\theta(u_{i,t} \mid \mathbf{f}_{i,t})
  }{
    \pi_{\theta_{\text{old}}}(u_{i,t} \mid \mathbf{f}_{i,t})
  }.
  \label{eq:is_ratio}
\end{equation}
With a shared team reward, we use a single centralized value function and a
shared GAE estimate~\cite{schulman2016gae}:
\begin{equation}
  \delta_t = r_t^{\text{team}} + \gamma\, V_\phi(\bar{\mathbf{h}}_{t+1})
  - V_\phi(\bar{\mathbf{h}}_t),\qquad
  \hat{A}_t = \sum_{l \ge 0}(\gamma \lambda_{\text{GAE}})^l \delta_{t+l},
  \label{eq:gae}
\end{equation}
and define the return target as
$R_t = \hat{A}_t + V_\phi(\bar{\mathbf{h}}_t)$.
The MAPPO loss is then
\begin{equation}
  \mathcal{L}_{\text{MAPPO}}
  = -\mathbb{E}_{t,i}\!\Big[\min\!\big(
      r_{i,t}(\theta)\,\hat{A}_t,\;
      \operatorname{clip}(r_{i,t}(\theta),\,1\!-\!\epsilon,\,
      1\!+\!\epsilon)\,\hat{A}_t
    \big)\Big]
  + c_v\,\mathbb{E}_{t}\!\big[(V_\phi(\bar{\mathbf{h}}_t) - R_t)^2\big]
  - c_e\,\mathbb{E}_{t,i}\!\big[
      \mathcal{H}[\pi_\theta(\cdot \mid \mathbf{f}_{i,t})]
    \big],
  \label{eq:ppo}
\end{equation}
where $\gamma\!=\!0.99$, $\lambda_{\text{GAE}}\!=\!0.95$,
$\epsilon\!=\!0.2$, $c_v\!=\!0.5$, $c_e\!=\!0.01$.
Rewards are normalized with running mean and variance.  Training uses a
three-stage curriculum: easy (3~UAVs, $50^2$ grid) $\to$ medium (6~UAVs,
$75^2$) $\to$ hard (10~UAVs, $100^2$), advancing when the 20-episode moving
average of search coverage exceeds a stage-specific threshold.
The full Stage~I training completes in ${\sim}10$\,h on a server with
$8\!\times\!$A100 GPUs.

% --- 4.4 TAGA-LLM Distillation ---
\subsection{Stage II: TAGA-LLM Distillation}
\label{sec:distill}

The trained TAGA-MAPPO policy achieves strong performance but remains a black
box.  Stage~II distills it into an LLM for interpretability and instruction
control.

\paragraph{Data collection.}
Let $\pi_E$ denote the converged TAGA-MAPPO expert.  We deploy $\pi_E$ for
3{,}000 episodes across three scenario configurations (varying target/threat
layouts).  From each episode~$e$ we uniformly sample one \emph{anchor scene}
per agent, yielding $N_{\text{anchor}} \!=\! 3{,}000$ distinct
observation--action pairs.  For each anchor scene with observation~$o_t$ and
directive~$d$, the expert action and training target are
\begin{equation}
  a_t^{*} = \pi_E(o_t,\, d),\qquad
  y_t = [\texttt{<think>},\, r_t,\, \texttt{</think>},\,
         \texttt{<act>},\, a_t^{*},\, \texttt{</act>}],
  \label{eq:teacher}
\end{equation}
where $r_t$ is a template-generated rationale that verbalizes the key
observation features and the rule triggered by the expert policy (not produced
by an LLM).  To improve instruction controllability, each anchor scene is
re-labeled under all five directives
(\textsc{bal}/\textsc{sch}/\textsc{atk}/\textsc{evd}/\textsc{coo}), producing
$N_{\text{cf}} = 5 \times N_{\text{anchor}} = 15{,}000$ counterfactual
samples.  The combined pool is
$N_{\text{raw}} = N_{\text{anchor}} + N_{\text{cf}} = 18{,}000$; after
filtering the top-75\% by episode return, the final dataset contains
$|\mathcal{D}| \approx 45$K training samples (some episodes contribute
multiple agents, and the 75\% filter is applied per-episode before
aggregation).

\paragraph{Prompt encoding.}
Each sample is a structured prompt comprising four blocks: (a)~self-state
(position, velocity, energy, mode), (b)~neighbor list sorted by ascending
distance, where each entry is annotated with TAGA-derived signal quality
(H/M/L, from joint thresholds on $(f_d, f_t)$) and numeric freshness ($f_t$),
(c)~visible target list with
coordinates and status, and (d)~an operational instruction
(\textsc{bal}/\textsc{sch}/\textsc{atk}/\textsc{evd}/\textsc{coo}).
Priority-based truncation (max 256 tokens) removes the most distant neighbors
first to stay within the context window.

\paragraph{Training.}
We fine-tune Qwen-4B~\cite{qwen2024} with LoRA~\cite{hu2022lora}
($r\!=\!16$, $\alpha\!=\!32$, dropout 0.05) on all attention/FFN projections
($\sim$6.3M trainable params).  Each sample follows the target format in
\eqref{eq:teacher}.  Let $\mathcal{M}_{\text{think}}$ and
$\mathcal{M}_{\text{act}}$ denote the sets of token positions within the
\texttt{<think>} and \texttt{<act>} spans, respectively.  The supervised
fine-tuning loss is a weighted sum of the two components:
\begin{equation}
  \mathcal{L}_{\text{SFT}}
  = -\frac{\lambda_{\text{act}}}{|\mathcal{M}_{\text{act}}|}
    \sum_{t \in \mathcal{M}_{\text{act}}} \log p_\theta(x_t \mid x_{<t})
    -\frac{\lambda_{\text{think}}}{|\mathcal{M}_{\text{think}}|}
    \sum_{t \in \mathcal{M}_{\text{think}}} \log p_\theta(x_t \mid x_{<t}),
  \label{eq:sft}
\end{equation}
where $\lambda_{\text{act}} \!=\! 1.0$ and
$\lambda_{\text{think}} \!=\! 0.5$, reflecting that action correctness is the
primary objective while the rationale serves as an auxiliary supervision signal
(the template-generated rationale provides structured reasoning cues but is not
a ground-truth explanation).
Training completes in $\sim$37\,min on a single A100 (from the same
$8\!\times\!$A100 server) with bf16 and gradient
checkpointing.  At inference, greedy decoding with 4-bit quantization achieves
$\sim$70\,ms per step.

The complete training procedure is summarized in Algorithm~\ref{alg:taga_llm}.

\begin{algorithm}[t]
\caption{Two-stage TAGA-LLM training and distillation procedure.}
\label{alg:taga_llm}
\small
\begin{algorithmic}[1]
\REQUIRE Environment config $\mathcal{C}$, pre-trained Qwen-4B, curriculum thresholds $\{\tau_1, \tau_2\}$
\ENSURE Distilled TAGA-LLM agent
\STATE \textit{Stage I --- TAGA-MAPPO training}
\STATE Initialize actor $\pi_\theta$, critic $V_\phi$, TAGA params $\lambda_d\!=\!1.0,\, \lambda_m\!=\!1.0,\, \lambda_t\!=\!0.1$
\STATE Set curriculum stage $\gets$ \textsc{easy} (3 UAVs, $50^2$ grid)
\FOR{update $= 1, \ldots, T_{\max}$}
  \STATE Collect rollout episodes with current $\pi_\theta$
  \STATE Encode observations $\to$ TAGA features $\tilde{\mathbf{h}}_i$ (Eq.~\ref{eq:taga_s}--\ref{eq:taga_agg})
  \STATE Compute actions via hierarchical policy (Sect.~\ref{sec:encoder_policy})
  \STATE Update $\theta, \phi$ via clipped PPO (Eq.~\ref{eq:ppo})
  \IF{20-episode moving avg.\ coverage $> \tau_{\text{stage}}$}
    \STATE Advance curriculum stage
  \ENDIF
\ENDFOR
\STATE \textit{Stage II --- LLM distillation}
\STATE Deploy converged $\pi_E \!=\! \pi_\theta^*$ for 3K episodes; sample $N_{\text{anchor}}\!=\!3$K anchor scenes
\STATE Re-label each anchor under 5 directives $\to$ $N_{\text{cf}}\!=\!15$K counterfactual samples
\STATE Filter top-75\% episodes $\to$ dataset $\mathcal{D}$ (${\approx}\,45$K samples)
\STATE Attach LoRA ($r\!=\!16$, $\alpha\!=\!32$) to Qwen-4B
\STATE Fine-tune on $\mathcal{D}$ with weighted SFT loss (Eq.~\ref{eq:sft}; $\lambda_{\text{act}}\!=\!1.0$, $\lambda_{\text{think}}\!=\!0.5$)
\RETURN TAGA-LLM agent
\end{algorithmic}
\end{algorithm}

% sec_experiments.tex — Section 5: Experiments and Simulation Results
% ICGNC 2026, Springer svproc template

\section{Experiments and Simulation Results}
\label{sec:experiments}

\subsection{Simulation Setup}

We simulate a partially observable grid battlefield with local $11\!\times\!11$
patches (obs radius $=5$) and shared global rewards as in Sect.~\ref{sec:problem}.
The default full-scale configuration follows
\texttt{configs/default.yaml}: a $100\!\times\!100$ grid, $10$~UAVs,
$5$~targets, and $7$~threat zones, with communication radius
$R_{\mathrm{comm}}=20$ and $\texttt{max\_steps}=200$ per episode.

\textbf{TAGA-MAPPO training.} We adopt three-stage curriculum learning:
\begin{itemize}\setlength{\itemsep}{0pt}\setlength{\parskip}{0pt}
  \item \texttt{stage\_1\_easy}\,: $3$~UAVs, $2$~targets, $0$~threats, $30\!\times\!30$;
  \item \texttt{stage\_2\_medium}\,: $6$~UAVs, $3$~targets, $3$~threats, $50\!\times\!50$;
  \item \texttt{stage\_3\_hard}\,: $10$~UAVs, $5$~targets, $7$~threats, $100\!\times\!100$.
\end{itemize}
\noindent MAPPO uses GAE($\lambda=0.95$), PPO clip
$\varepsilon=0.2$, discount $\gamma=0.99$, Adam learning rate $3\times10^{-4}$,
and entropy coefficient $0.05$ for the short sanity runs reported below.

\textbf{Baselines for the quick validation run.}  We compare (i)~a \emph{random
joint policy} that uniformly samples hierarchical actions on the
\texttt{stage\_3\_hard} layout, and (ii)~a \emph{greedy TAGA-MAPPO policy}
(argmax over mode and move logits) after a \emph{limited} training budget of
$90$~PPO updates on CPU, intended as an engineering sanity check rather than a
fully converged benchmark.  Full-scale comparisons against classical planners
(e.g., ACO), flat MAPPO, and TAGA-LLM under five directives
(\textsc{bal}, \textsc{sch}, \textsc{atk}, \textsc{evd}, \textsc{coo}) are left
to the extended evaluation track with $30$~seeds and Welch's $t$-test at
$\alpha=0.05$, as originally planned.

\textbf{Hardware:} $8\!\times\!$A100 GPU server; Stage~I training ${\sim}10$\,h,
Stage~II LoRA fine-tuning ${\sim}37$\,min (single GPU).

\subsection{Simulation Results}

\begin{table}[t]
\centering
\caption{Rapid local validation (mean $\pm$ std).  Random: 24 episodes on
\texttt{stage\_3\_hard}.  TAGA-MAPPO: greedy rollouts after 90 updates
(10 episodes on the post-training environment layout).}
\label{tab:sanity}
\begin{tabular}{lcc}
\toprule
Metric & Random & TAGA-MAPPO (greedy, short train) \\
\midrule
Search coverage      & $0.361 \pm 0.043$ & $0.557 \pm 0.068$ \\
Target destroy rate  & $0.183 \pm 0.143$ & $0.000 \pm 0.000$ \\
Norm.\ completion time & $0.900 \pm 0.000$ & $1.000 \pm 0.000$ \\
UAV survival rate    & $0.742 \pm 0.125$ & $0.620 \pm 0.132$ \\
Coop.\ attack ratio  & $0.000 \pm 0.000$ & $0.000 \pm 0.000$ \\
\bottomrule
\end{tabular}
\end{table}

Table~\ref{tab:sanity} shows that even with a small optimization budget,
TAGA-MAPPO already improves search coverage relative to random exploration;
the destroy-rate and time-to-completion metrics indicate that substantially
longer training (and LLM distillation) are required before claiming mission
completion at scale---consistent with the staged experimental roadmap.

\subsection{Instruction Controllability}

We will evaluate TAGA-LLM under five operational directives
(\textsc{bal}, \textsc{sch}, \textsc{atk}, \textsc{evd}, \textsc{coo}) using
the topology-aware prompt encoder described in Sect.~\ref{sec:distillation}.
The distilled model is expected to shift coverage--aggression trade-offs
measurably across directives; quantitative curves are omitted here pending
completion of the LoRA fine-tuning pipeline.

\subsection{Ablation Studies}

Planned ablations include (i)~TAGA-aware prompt encoding versus a flat text
serialization; (ii)~chain-of-thought supervision versus action-only supervision;
(iii)~counterfactual data augmentation on/off; (iv)~ablating individual TAGA
factors ($\lambda_d, \lambda_m, \lambda_t$) in the expert policy.  Each
ablation will be reported with the same five metrics as the main table once the
long-run training budget is finalized.

\subsection{Scalability and Robustness}

We target scalability sweeps with $5/10/15/20$~UAVs, communication dropout at
$10\%/20\%/30\%$, and inference latency benchmarks for the distilled LLM.
These experiments are not included in the rapid CPU sanity run above.

% TODO: 图片文件缺失，暂时注释
% \begin{figure}[t]
% \centering
% \includegraphics[width=\textwidth]{training_curves.pdf}
% \caption{TAGA-MAPPO training traces under curriculum scheduling
% (policy/value loss, entropy, reward/coverage).  Vertical markers indicate
% automatic stage transitions when rolling coverage exceeds stage thresholds.}
% \label{fig:training_curves}
% \end{figure}
%
% \begin{figure}[t]
% \centering
% \includegraphics[width=0.6\textwidth]{curriculum_stage_timeline.pdf}
% \caption{Curriculum stage occupancy over PPO updates for the same sanity run.}
% \label{fig:curriculum_timeline}
% \end{figure}

% sec_conclusion.tex — Section 6: Conclusion (compressed)
% ICGNC 2026, Springer svproc template

\section{Conclusion}
\label{sec:conclusion}

This paper presented TAGA-LLM, a two-stage framework that bridges MARL cooperative
performance with LLM interpretability for multi-UAV mission planning.  The TAGA
mechanism introduces three learnable, physics-informed attention
factors---distance decay $f_d$, task awareness $f_m$, and information freshness
$f_t$---that model realistic communication constraints beyond static-graph
assumptions.  The MARL-to-LLM distillation pipeline, combining topology-aware
prompt encoding, chain-of-thought supervision, and counterfactual augmentation,
transfers cooperative strategies into a Qwen-4B model that provides transparent
reasoning, instruction controllability under the five mission directives
(\textsc{bal}--\textsc{coo}), and efficient edge deployment
(${\sim}70$\,ms/step with 4-bit quantization).

\paragraph{Limitations.}
The current framework assumes homogeneous UAVs and a simplified grid-based
environment.  The distillation quality depends on the coverage of the expert
trajectory dataset, and the 256-token prompt truncation may lose information in
dense scenarios.  Additionally, the chain-of-thought supervision relies on
template-based rationale generation rather than human annotations.

\paragraph{Future work.}
Promising directions include extending to heterogeneous swarms with
role-specific policies, incorporating online adaptation for non-stationary
environments, and validating on real-world UAV testbeds with physical
communication constraints.


% ============================
% 参考文献
% ============================
\bibliographystyle{styles/bibtex/splncs03_unsrt}
\bibliography{references}

\end{document}
